\documentclass[11pt,a4paper]{article}
\usepackage{amsmath,amsfonts,amssymb,amsthm}
\usepackage[OT1, OT2]{fontenc}
\usepackage[english,serbian]{babel}
\usepackage{paralist}
\usepackage{algorithm}
\usepackage{titlesec}

\topmargin       -11.0mm
\oddsidemargin      0mm
\evensidemargin     0mm
\textheight     235.5mm
\textwidth      161.0mm
\columnsep        4.1mm
\parindent        1.0em
\headsep          6.3mm
\headheight        12pt
\lineskip           1pt
\normallineskip     1pt
\def\baselinestretch{1.0}

\newcommand{\lN}{\mathbb{N}}
\newcommand{\lR}{\mathbb{R}}

\newtheorem{thm}{Teorema}[section]
\renewcommand{\thethm}{\arabic{section}.\arabic{thm}}

\newtheorem{tvr}[thm]{Tvrdjenje}

\newtheorem{posl}[thm]{Posledica}

\newtheorem{defi}{Definicija}[section]
\renewcommand{\thedefi}{\arabic{section}.\arabic{defi}}

\newtheorem*{nap}{Napomena}

\newtheorem{lem}{Lema}[subsection]

\newtheorem{pri}{Primer}[section]
\renewcommand{\thepri}{\arabic{section}.\arabic{pri}}

\newtheorem{zad}{Zadatak}

\newtheorem*{dokaz}{Dokaz}
%\renewcommand{\thedokaz}{}

\newcommand{\cyr}{\fontencoding{OT2}\selectfont\selectlanguage{serbian}}
\newcommand{\lat}{\fontencoding{T1}\selectfont\selectlanguage{english}}

\newtheorem{assumption}{Assumption}[section]
\renewcommand{\theassumption}{\arabic{section}.\arabic{assumption}}

\newtheorem{theorem}{Theorem}[section]
\renewcommand{\thetheorem}{\arabic{section}.\arabic{theorem}}

\newtheorem{lemma}{Lemma}[section]
\renewcommand{\thelemma}{\arabic{section}.\arabic{lemma}}

\newtheorem{corollary}{Corollary}[section]
\renewcommand{\thecorollary}{\arabic{section}.\arabic{corollary}}

\newtheorem{proposition}{Proposition}[section]
\renewcommand{\theproposition}{\arabic{section}.\arabic{proposition}}

\newcounter{appendixcounter}

\newtheorem{theoremappendix}{Theorem}
\renewcommand{\thetheoremappendix}{A\arabic{appendixcounter}.\arabic{theoremappendix}}

%\newtheorem{lemmaappendix}{Lemma}
%\renewcommand{\thelemmaappendix}{A\arabic{appendixcounter}.\arabic{lemmaappendix}}

\newtheorem{lemmaappendix}{Lemma}
\renewcommand{\thelemmaappendix}{A.\arabic{lemmaappendix}}

\newtheorem{assumptionappendix}{Assumption}[appendixcounter]
\renewcommand{\theassumptionappendix}{A\arabic{appendixcounter}.\arabic{assumptionappendix}}


\newtheorem*{remark}{Remark}%[section]
%\renewcommand{\theremark}{\arabic{section}.\arabic{remark}} 


\newtheorem{step}{Step}%[section]
\renewcommand{\thestep}{\arabic{step}} 

\titleformat{\section}[hang]{\bf\Large}{\thesection.}{0.4em}{}
\titleformat{\subsection}[hang]{\bf\large}{\thesubsection.}{0.4em}{}

\renewcommand{\baselinestretch}{1.01}

\begin{document}

\title{Konvergencija nizova} 

\author{Milica V. Tadiћ
\thanks{Matamatichka Gimnazija, Kraljice Natalije 37, Beograd  }}

\date{}

\maketitle


\section{Uvod} \label{section0}

 U ovom seminarskom radu razmatramo nizove u skupu realnih brojeva i njihove osnovne asimpto\cyr{t}\cyr{s}ke osobine. 
 Fokus rada je na konceptu konvergencije i 
 granichne vrednosti niza realnih brojeva. 
 Seminarski rad obuhvata osnovne osobine kovergentnih nizova i njihovih granichnih vrednosti. 
 Veћина особина (razmatranih u radu) 
 potkrepljena je dokazima i ilustrovana primerima. 
 
 Seminarski rad je organizivan na sledeћи nachin. 
 U odeljku \ref{section1}, definishemo  nizove u skupu realnih brojeva i razmatramo njihove osnovne algebarske osobine. 
 Koncept konvergencije i granichne vrednosti niza realnih brojeva definishemo u odeljku \ref{section2}. 
 Srodni koncepti (kao shto su odredjeno divergentni i nula-nizovi) takodje se definishu u odeljku \ref{section2}. 
 %U odeljaku \ref{section3}, razmaramo osnovna 
 %asimpto\cyr{t}\cyr{s}ka svojstva nizova realnih brojeva. 
 %U odeljku \ref{section4}, 
%data su dva (netrivijana) primera primene rezultata razmatranih u 
%odeljcima \ref{section2} i \ref{section3}. 
 
\section{Nizovi u skupu realnih brojeva} \label{section1}

Niz realnih brojeva je jedan od najvazhnijih pojmova 
u matematichkoj analizi i algebri. 
Niz realnih brojeva mozhemo zamisliti kao vektor sa beskonachno 
mnogo elemenata, 
odnosno kao pravilo koje svakom prirodnom broju $n$
pridzhuje realan broj $a_n$. 
U tom sluchaju, realan broj $a_n$ prestavlja $n$-ti chlan niza, 
a niz prikazujemo ovako:
$(a_n)_{n\in\mathbb{N}}$ 
(ili $\{a_n\}_{n\geq 1}$). 
Niz realnih brojeva i njegovi chlanovi se definishu 
na sledeћi nachin: 

\begin{defi}
\textbf{Niz realnih brojeva} je preslikavanje skupa $\mathbb{N}$  
u skup $\mathbb{R}$. 
Vrednost ovog preslikavanja za $n\in\mathbb{N}$ je $n$-ti 
chlan niza. 
\end{defi}

U matematichoj analizi, znachajnu ulogu imaju ogranicheni i 
monotoni nizovi. Ove nizove definishemo na 
sledeћi nachin: 

\begin{defi}
Neka je $(a_n)_{n\in\mathbb{N}}$ niz realnih brojeva.  

(a) $(a_n)_{n\in\mathbb{N}}$ je \textbf{ogranichen odozgo} ako postoji realan broj $b$ takav da je $a_n\leq b $ za sve $n \in \lN$.

(b) $(a_n)_{n\in\mathbb{N}}$ je \textbf{ogranichen odozdo} ako postoji realan broj $c$ takav da je $a_n \geq c $ za sve $n \in \lN$.

(v) $(a_n)_{n\in\mathbb{N}}$ je \textbf{ogranichen} ako je ogranichen odozdo i odozgo, tj. ako postoje realni brojevi $b$ i $c$ takvi da je $c \leq a_n \leq b$ za svako $n \in \lN$.
\end{defi}

\begin{defi}
Neka je $(a_n)_{n\in\mathbb{N}}$ niz realnih brojeva.  

(a) $(a_n)_{n\in\mathbb{N}}$ je 
\textbf{rastuћи} ako je $a_n \leq a_{n+1}$
za svako $n\in\mathbb{N}$. 

(b)  
Ako je $a_n < a_{n+1}$
za svako $n\in\mathbb{N}$, onda je 
$(a_n)_{n\in\mathbb{N}}$ 
\textbf{strogo rastuћи}. 

(v) $(a_n)_{n\in\mathbb{N}}$ je 
\textbf{opadajuћи} ako je $a_n \geq a_{n+1}$
za svako $n\in\mathbb{N}$. 

(g)  
Ako je $a_n > a_{n+1}$
za svako $n\in\mathbb{N}$, onda je 
$(a_n)_{n\in\mathbb{N}}$ 
\textbf{strogo opadajuћи}. 

(dj)
Ako $(a_n)_{n\in\mathbb{N}}$ zadovoljava bilo koji od uslova
(a) -- (g), onda je $(a_n)_{n\in\mathbb{N}}$ \textbf{monoton}. 

(e) $(a_n)_{n\in\mathbb{N}}$ je \textbf{strogo monoton}
ako $(a_n)_{n\in\mathbb{N}}$ zadovoljava uslov (b) ili (d). 
\end{defi}

Prethodne definicije ilustrovane su sledeћim primerima. 

\begin{pri}
Neka je $(a_n)_{n\in\mathbb{N}}$ niz realnih brojeva. 

(a) $(a_n)_{n\in\mathbb{N}}$ je {\bf aritmetichki niz} 
ako je $a_{n+1}=a_{n}+d$ za neko $d\in\mathbb{R}$ i za svako 
$n\in\mathbb{N}$. 
U tom sluchaju, imamo $a_n=a_1+(n-1)d$ za svako $n\in\mathbb{N}$.  
Iz ovog sledi da je $(a_n)_{n\in\mathbb{N}}$ monotono strogo 
rastuћi i ogranichen odozdo za $d>0$, kao i da je 
$(a_n)_{n\in\mathbb{N}}$ monotono strogo 
opadajuћi i ogranichen odozgo za $d<0$. 

(b) $(a_n)_{n\in\mathbb{N}}$ je {\bf geometrijski niz} 
ako je $a_{n+1}=a_{n}q$ za neko $q\in\mathbb{R}$ i za svako 
$n\in\mathbb{N}$. 
U tom sluchaju, imamo $a_n=a_1q^{n-1}$ za svako $n\in\mathbb{N}$.  
Iz ovog sledi da je $(a_n)_{n\in\mathbb{N}}$ monotono strogo 
rastuћi i ogranichen za $a_1<0$, $0<q<1$, kao i da je 
$(a_n)_{n\in\mathbb{N}}$ monotono strogo 
opadajuћi i ogranichen za $a_0>0$, $0<q<1$.
Mozhe se takodje pokazati da je $(a_n)_{n\in\mathbb{N}}$ je ogranichen 
za $|q|<1$. 

(v) $(a_n)_{n\in\mathbb{N}}$ je {\bf harmonijski niz} 
ako je $a_1=1$, $a_2=1/2$ i $a_{n+2}=2/(1/a_{n+1}+1/a_n)$ za svako 
$n\in\mathbb{N}$. 
U tom sluchaju, imamo $a_n=1/n$ za svako $n\in\mathbb{N}$.  
Iz ovog sledi da je $(a_n)_{n\in\mathbb{N}}$ monotono strogo 
opadajuћi i ogranichen. 
\end{pri}

Slichno kao u skupu realnih brojeva, 
moguce je definisati operacije sabiranja, oduzimanja, mnozhenja i 
deljenja za nizove relnih brojeva. 

\begin{defi}
Neka su $(a_n)_{n\in\mathbb{N}}$ i $(b_n)_{n\in\mathbb{N}}$
nizovi realnih brojeva. 

(a) $(a_n)+(b_n)$ je niz chiji je $n$-ti chlan 
$a_n+b_n$. 

(b) $(a_n)-(b_n)$ je niz chiji je $n$-ti chlan 
$a_n-b_n$. 

(v) $(a_n)\cdot(b_n)$ je niz chiji je $n$-ti chlan 
$a_nb_n$. 

(g) $(a_n)/(b_n)$ je niz chiji je $n$-ti chlan 
$a_n/b_n$ (pod pretpostavkom da je $b_n\neq 0$). 

(d) $c\cdot(a_n)$ je niz chiji je $n$-ti chlan 
$ca_n$, gde je $c\in\mathbb{R}$. 
\end{defi}

\section{Konvergencija nizova i srodni pojmovi} \label{section2}

U problemima matematichke analize,
chesto je potrebno utvrditi da li postoji zakonitost 
u ponashanju niza realnih brojeva kada indeks niza
tezhi ka beskonachnosti. 
Ispoljeno ponashanje niza realnih brojeva kada indeks tezhi  
beskonachnosti naziva se {\bf asimpto\cyr{t}\cyr{s}kim ponashanjem}, 
dok se odgovarajuћe osobine zovu 
{\bf asimpto\cyr{t}\cyr{s}ke osobine}. 
Konvergencija i granichna vrednost su najvazhnije 
asimpto\cyr{t}\cyr{s}ke 
osobine u matematichkoj analizi. 
Da bih se definisale ove osobine, koristimo 
koncept $\varepsilon$-okoline. 
$\varepsilon$-okolina se definishe na sledeћi nachin: 

 \begin{defi}
 Neka su $a$ i $\varepsilon$ realni brojevi i neka je $\varepsilon>0$. 
 Onda je \textbf{$\varepsilon$-okolina broja $a$} interval
 \begin{align*}
     U_{\varepsilon}(a) 
     =
     (a-\varepsilon,a+\varepsilon)
     =
     \{x\in\mathbb{R}: |x-a|<\varepsilon\}. 
 \end{align*}
 \end{defi}
 
 Koristeћi $\varepsilon$-okolinu, definishemo 
 granichnu vrednost niza. 
 
 \begin{defi}
 Neka je $(a_n)_{n\in\mathbb{N}}$ niz realnih brojeva. 
 Realni broj $a$ je onda \textbf{granichna vrednost (ili limes)} 
 niza $(a_n)_{n \in \lN}$ ako za svaki pozitivan realan broj $\varepsilon$ postoji prirodan broj $n_0(\varepsilon)$ takav da 
 $a_n\in U_{\varepsilon}(a)$ za svako $n>n_{0}(\varepsilon)$, 
 tj, ako vazhi 
 \begin{align*}
     (\forall\varepsilon\in(0,\infty) )
     \; (\exists n_0(\varepsilon)\in\mathbb{N} )
     \; (\forall n\in\mathbb{N})
     \; (n>n_0(\varepsilon)\Rightarrow
     |a_n-a|<\varepsilon).  
 \end{align*}
 Za ovo upotrebljavamo notaciju 
 $a_n\rightarrow a$ ili $\displaystyle\lim_{n\rightarrow\infty}a_n=a$, 
 i kazhemo da $(a_n)_{n\in\mathbb{N}}$ tezhi (konvergira) 
 ka $a$. 
 \end{defi}
 
 Na osnovu prethodne definicije, imamo da je 
 $\displaystyle\lim_{n\rightarrow\infty}a_n=a$ ako i samo ako 
 za svaki (proizvoljno mali) 
 pozitivan realan broj $\varepsilon$, 
 $\epsilon$-okolina broja $a$ sadrzhi sve 
 chlanove niza $(a_n)_{n\in\mathbb{N}}$, 
 osim njih konachno mnogo. 
 Kako $\varepsilon$-okolina broja $a$ postaje sve ,,uzha i uzha'' 
 kada se $\varepsilon$ smanjuje i priblizhava nuli, 
 zakljuchujemo da su $a_n$ i 
 $a$ priblizhno isti  
 za svaki ,,dovoljno veliki indeks $n$'', 
 tj. $a_n\approx a$ za $n\gg 1$. 
 To znachi da se chlanovi niza $(a_n)_{n\in\mathbb{N}}$ 
 ,,grupishu oko $n$'' kada $n$ tezhi beskonachnosti. 
 
 Koristeћi granichnu vrednost, definishemo konvergentan,  
divergentan i nula niz. 
 
\begin{defi}
(a) Niz realnih brojeva je \textbf{konvergentan} ukoliko ima granchnu vrednost. 

(b) Niz realnih brojeva je \textbf{divergentan} ako 
nema granichnu vrednost.

(v) \textbf{Nula-niz} je niz realnih brojeva chija granichna 
vrednost je nula. 
\end{defi}

Prethodne definicije ilustrovane su sledeћim primerima

\begin{pri}\label{pri3.1}
(a) Neka $a=1$ i neka je $(a_n)_{n\in\mathbb{N}}$ niz chiji 
je $n$-ti chlan $a_n=(n-1)/n$. 
Lako se onda pokazuje da je $|a_n-a|=1/n$. 
Neka je $\varepsilon$ bilo koji pozitivan realan broj, 
dok je $n_{0}(\varepsilon)$ najmanji prirodan broj 
veћi od $1/\varepsilon$ 
(tj. $n_{0}(\varepsilon)=\lceil 1/\varepsilon \rceil$).
Kako je $n_{0}(\varepsilon)\geq 1/\varepsilon$, dobijamo 
\begin{align*}
    n>n_{0}(\varepsilon)
    \;\Rightarrow\; 
    n>1/\varepsilon
    \;\Rightarrow\;  
    |a_n-a|=1/n<\varepsilon. 
\end{align*}
Koristeћi onda definiciju granichne vrednosti 
i konvergentnog niza, 
zakljuchujemo da je 
$\displaystyle\lim_{n\rightarrow\infty}a_n=a=1$ i da je 
$(a_n)_{n\in\mathbb{N}}$ konvergentan. 

(b)  Neka je $a=0$ i neka je $(a_n)_{n\in\mathbb{N}}$ niz chiji 
je $n$-ti chlan $a_n=(-1)^n/n^2$. 
Lako se onda pokazuje da je $|a_n-a|=1/n^2$. 
Neka je $\varepsilon$ bilo koji pozitivan realan broj, 
dok je $n_{0}(\varepsilon)$ najmanji prirodan broj 
veћi od $1/\sqrt{\varepsilon}$ 
(tj. $n_{0}(\varepsilon)=\lceil 1/\sqrt{\varepsilon} \rceil$).
Kako je $n_{0}(\varepsilon)\geq 1/\sqrt{\varepsilon}$, dobijamo 
\begin{align*}
    n>n_{0}(\varepsilon)
    \;\Rightarrow\; 
    n>1/\sqrt{\varepsilon}
    \;\Rightarrow\;  
    |a_n-a|=1/n^2<\varepsilon. 
\end{align*}
Koristeћi onda definiciju granichne vrednosti 
i nula-niza, 
zakljuchujemo da je 
$\displaystyle\lim_{n\rightarrow\infty}a_n=a=0$ i da je
$(a_n)_{n\in\mathbb{N}}$ nula niz. 

(v)  Neka je $(a_n)_{n\in\mathbb{N}}$ niz chiji 
je $n$-ti chlan $a_n=(-1)^n$. Pokazhimo da je $(a_n)_{n\in\mathbb{N}}$
divergentan. 
Pretpostavimo suprotno. 
Onda postoji realan broj $a$ i za svaki pozitivan realan broj $\varepsilon$
postoji prirodan broj $n_{0}(\varepsilon)$ tako da je 
$|a_n-a|<\varepsilon$ za $n>n_{0}(\varepsilon)$. 
U tom sluchaju dobijamo 
\begin{align*}
    &
    a=(a-a_{2n+1} ) + a_{2n+1}
    =-1+(a-a_{2n+1} )
    \leq 
    -1+|a-a_{2n+1} |
    \leq 
    -1/2, 
    \\
    &
    a=(a-a_{2n} ) + a_{2n}
    =1+(a-a_{2n} )
    \geq 
    1-|a-a_{2n} |
    \geq 
    1/2
\end{align*}
za $n>n_0(1/2)/2$. 
Kako prethodne dve nejednakosti ne mogu da vazhe istovremeno, 
zaklju\cyr{ch}ujemo da je $(a_n)_{n\in\mathbb{N}}$ divergentan. 
\end{pri}

Osim konvergentnih nizova, 
u matematichkoj analizi se razmatraju i nizovi 
koji ,,odredjeno divergiraju'', tj tezhe (konvergiraju) ka 
$+\infty$ ili $-\infty$. 
Ovi se nizovi definishu na sledeћi nachin. 

\begin{defi}
Neka je $(a_n)_{n\in\mathbb{N}}$ niz realnih brojeva. 

(a) $(a_n)_{n\in\mathbb{N}}$ konvergira ka $+\infty$
ako za svaki pozitivni realni broj $M$ postoji prirodan broj
$n_{0}(M)$ tako da je 
$a_{n}>M$ za svako $n>n_{0}(M)$, 
tj. ako vazhi 
\begin{align*}
     (\forall M\in(0,\infty) )
     \; (\exists n_0(M)\in\mathbb{N} )
     \; (\forall n\in\mathbb{N})
     \; (n>n_0(M)\Rightarrow
     a_n>M).  
 \end{align*}
 Za ovo upotrebljavamo notaciju 
 $a_n\rightarrow \infty$ ili
 $\displaystyle\lim_{n\rightarrow\infty}a_n=\infty$. 
 
(b) $(a_n)_{n\in\mathbb{N}}$ konvergira ka $-\infty$
ako za svaki negativni realni broj $M$ postoji prirodan broj
$n_{0}(M)$ tako da je 
$a_{n}<M$ za svako $n>n_{0}(M)$, 
tj. ako vazhi 
\begin{align*}
     (\forall M\in(-\infty,0) )
     \; (\exists n_0(M)\in\mathbb{N} )
     \; (\forall n\in\mathbb{N})
     \; (n>n_0(M)\Rightarrow
     a_n<M).  
 \end{align*}
 Za ovo upotrebljavamo notaciju 
 $a_n\rightarrow -\infty$ ili
 $\displaystyle\lim_{n\rightarrow\infty}a_n=-\infty$.
\end{defi}

Na osnovu prethodne definicije, 
imamo da je $\displaystyle\lim_{n\rightarrow\infty}a_n=\infty$
ako i samo ako je svaki (proizvoljno veliki) pozitivan broj $M$
manji od svih chlanova niza $(a_n)_{n\in\mathbb{N}}$, 
osim od konachno mnogo njih.  
Kako je $M$ proizvoljno velika pozitivna velichina, 
zakljuchujemo da je i $a_n$ proizvoljno velika vechina 
za svaki ,,dovoljno veliki indeks $n$'', 
tj. $a_n\gg 1$ za $n\gg 1$. 
Na slichan nachin mozhemo protumachiti asimpto\cyr{t}\cyr{s}ku 
osobinu 
$\displaystyle\lim_{n\rightarrow\infty}a_n=-\infty$. 


Sledeћi primer ilustruje odredjeno divergentne nizove. 

\begin{pri}
(a)
Neka je $(a_n)_{n\in\mathbb{N}}$ niz chiji 
je $n$-ti chlan $a_n=n^2$. 
Pored toga, neka je $M$ bilo koji pozitivan realan broj, 
dok je $n_{0}(M)$ najmanji prirodan broj veћi od 
$\sqrt{M}$ (tj. $n_{0}(M)=\lceil \sqrt{M} \rceil$). 
Kako je $n_{0}(M)\geq \sqrt{M}$, dobijamo 
\begin{align*}
    n\geq n_{0}(M)
    \;\Rightarrow\; 
    n>\sqrt{M} 
    \;\Rightarrow\; 
    a_n=n^2> M. 
\end{align*}
Koristeћi onda definiciju odredjeno divergentnog niza, 
zakljuchujemo da je 
$\displaystyle\lim_{n\rightarrow\infty}a_{n}=\infty$. 

(b) Neka su $a$ i $d$ realni brojevi takvi da je 
$d>0$ i neka je 
$(a_n)_{n\in\mathbb{N}}$ aritmetichki niz chiji je 
$n$-ti chlan $a_{n}=a+(n-1)d$. 
Pored toga, neka je $M$ bilo koji pozitivan realan broj, 
dok je $n_{0}(M)$ najmanji prirodan broj 
veћi od $1+(|a|+M)/d$
(tj., $n_0(M)=\lceil 1+(|a|+M)/d \rceil$). 
Kako je $n_{0}(M)\geq 1+(|a|+M)/d$, dobijamo 
\begin{align*}
    n\geq n_{0}(M)
    \;\Rightarrow\; 
    n>1+(|a|+M)/d
    %\;\Rightarrow\; 
    %(n-1)d-|a|>M
    \;\Rightarrow\; 
    a_n=a+(n-1)d\geq (n-1)d-|a|>M. 
\end{align*}
Koristeћi onda definiciju odredjeno divergentnog niza, 
zakljuchujemo da je 
$\displaystyle\lim_{n\rightarrow\infty}a_{n}=\infty$. 
\end{pri}


\section{Asimpto\cyr{t}\cyr{s}ke osobine konvergentnih 
i odredjeno divergentnih nizova} \label{section3}

U ovom odeljku razmatramo asimpto\cyr{t}\cyr{s}ko ponashanje 
konvergentnih i odredjeno divergentnih nizova, 
kao i njihove osnovne osobine. 

\begin{thm}\label{thm1}
Neka je $(a_n)_{n\in\mathbb{N}}$ konvergentan niz realnih brojeva. 
Onda je njegova granichna vrednost jedinstvena. 
\end{thm}

\begin{nap}
Teorema \ref{thm1} pokazuje da nije moguћe da niz realnih brojeva 
ima dve ili vishe granichnih vrednosti. 
\end{nap}

\begin{dokaz}
Pretpostavimo suprotno. 
Onda niz $(a_n)_{n\in\mathbb{N}}$ ima dve razlichite granichne 
vrednosti $a$ i $b$. 
Koristeћi definiciju granichne vrednosti, 
zakljuchujemo da za svaki pozitivan realan broj $\varepsilon$ 
postoje prirodni brojevi $k_{0}(\varepsilon)$ i $l_{0}(\varepsilon)$ 
tako da je 
\begin{align}\label{t1.1}
    n>k_{0}(\varepsilon) \;\Rightarrow\; |a_{n}-a|<\varepsilon, 
    \;\;\;\;\; 
    n>l_{0}(\varepsilon) \;\Rightarrow\; |a_n-b|<\varepsilon. 
\end{align}
Neka je $c=|a-b|$. Kako je $a\neq b$, imamo $c>0$. 
Onda (\ref{t1.1}) daje 
\begin{align}\label{t1.3}
    c=|a-b|=|(a_n-a)-(a_n-b)|
    \leq 
    |a_n-a|+|a_n-b|
    \leq 
    \frac{2c}{3}
\end{align}
za $n>\max\{k_{0}(c/3),l_{0}(c/3)\}$.
Iz (\ref{t1.3}) sledi $c/3<0$. 
Ovo je nemoguћe, jer je $c>0$. 
\end{dokaz}

\begin{thm}\label{thm2}
Neka je $(a_n)_{n\in\mathbb{N}}$ konvergentan niz realnih brojeva
i neka je $a=\displaystyle\lim_{n\rightarrow\infty}a_n$.
Pored toga, neka je $(b_n)_{n\in\mathbb{N}}$ niz realnih brojeva 
chiji je $n$-ti chlan $b_{n}=a_n-a$. 
Onda imamo sledeћe: 

(a) $(a_n)_{n\in\mathbb{N}}$ je ogranichen niz. 

(b) $(b_n)_{n\in\mathbb{N}}$ je nula-niz. 
\end{thm}

\begin{nap}
Deo (a) teorema \ref{thm2} pokazuje da neogranichen niz 
nikad nije konvergentan. 
Deo (b) iste teoreme uspostvlja vezu izmedju konvergentnih i 
nula nizovam, tj. pokazuje 
da je $\displaystyle\lim_{n\rightarrow\infty}a_n=a$ 
ako i samo ako je $(a_n-a)_{n\in\mathbb{N}}$ nula-niz. 
\end{nap}

\begin{dokaz}
Na osnovu definicije granichne vrednosti, 
za svaki pozitivan realn broj $\varepsilon$ postoji prirodan broj 
$n_{0}(\varepsilon)$ tako da je 
\begin{align}\label{t2.1}
    n>n_{0}(\varepsilon)
    \;\Rightarrow\; 
    |a_n-a|<\varepsilon. 
\end{align}

(a)
Neka je $m=n_0(1)$ i neka je 
\begin{align*}
    c=\max\{|a_1|,\dots,|a_m|\}, 
    \;\;\;\;\; 
    d=\max\{|a|+1,c\}. 
\end{align*}
Onda, koristeћi (\ref{t2.1}), dobijamo 
\begin{align}\label{t2.3}
    |a_n|=|(a_n-a)+a|
    \leq 
    |a_n-a|+|a|
    \leq 
    1+|a|
    \leq 
    d
\end{align}
za $n>m=n_0(1)$. 
Pored toga, imamo  
\begin{align}\label{t2.5}
    |a_n|
    \leq 
    \max\{|a_1|,\dots,|a_m|\}
    =
    c
    \leq d
\end{align}
za $1\leq n\leq m$. 
Iz (\ref{t2.3}) i (\ref{t2.5}) sledi 
$|a_n|\leq d$ za $n\geq 1$
(tj. $-d\leq a_n\leq d$ za $n\in\mathbb{N}$).
Onda, na osnovu definicije ogranichenog niza, 
zakljuchujemo da je $(a_n)_{n\in\mathbb{N}}$ je obranichen. 

(b)
Neka je $b=0$. Onda (\ref{t2.1}) daje 
$|b_n-b|=|a_n-a|<\varepsilon$
za $n>n_0(\varepsilon)$ i svaki pozitivan realni broj $\varepsilon$. 
Koristeћi definiciju granichne vrednosti i nula niza 
zakljuchujemo onda da je 
$\displaystyle\lim_{n\rightarrow\infty}b_n=b=0$ i 
da je $(b_n)_{n\in\mathbb{N}}$ nula-niz. 
\end{dokaz}

\begin{thm}\label{thm3}
Neka su $(a_n)_{n\in\mathbb{N}}$ i $(b_n)_{n\in\mathbb{N}}$ konvergentni nizovi realnih brojeva. 
Onda imamo sledeћe: 

(a) Za sve realne brojeve $\alpha$ i $\beta$, 
niz $\alpha\cdot(a_n)_{n\in\mathbb{N}}+\beta\cdot(b_n)_{n\in\mathbb{N}}$
je konvergentan i ispunjava 
\begin{align*}
    \lim_{n\rightarrow\infty}
    (\alpha a_{n} + \beta b_{n} )
    =
    \alpha \lim_{n\rightarrow\infty} a_{n}
    + 
    \beta \lim_{n\rightarrow\infty} b_{n}. 
\end{align*}

(b) Niz $(a_n)_{n\in\mathbb{N}}\cdot(b_n)_{n\in\mathbb{N}}$
je konvergentan i ispunjava  
\begin{align*}
    \lim_{n\rightarrow\infty}
    (a_{n}b_{n} )
    =
    \lim_{n\rightarrow\infty} a_{n}
    \cdot
    \lim_{n\rightarrow\infty} b_{n}. 
\end{align*}

(v) Ako je $\displaystyle\lim_{n\rightarrow\infty}b_n\neq 0$ i ako je 
$b_n\neq 0$ za svako $n\in\mathbb{N}$, 
niz $(a_n)_{n\in\mathbb{N}}/(b_n)_{n\in\mathbb{N}}$
konvergentan i ispunjava  
\begin{align*}
    \lim_{n\rightarrow\infty}
    \frac{a_{n}}{b_{n} }
    =
    \frac{\displaystyle\lim_{n\rightarrow\infty} a_{n} }
    {\displaystyle\lim_{n\rightarrow\infty} b_{n} }. 
\end{align*}
\end{thm}

\begin{nap}
Teorema \ref{thm3} pokazuje kako se mozhe izrachunati 
granichna vrednost 
zbira, razlike, proizvoda i kolichnika konvergentnih nizova. 
\end{nap}

\begin{dokaz}
Neka je $a=\displaystyle\lim_{n\rightarrow\infty}a_n$ i 
$b=\displaystyle\lim_{n\rightarrow\infty}b_n=b$. 
Onda, na osnovu definicije granichne vrednosti, 
za svaki pozitivan realan broj $\varepsilon$ postoje 
prirodni brojevi $k_{0}(\varepsilon)$ i $l_{0}(\varepsilon)$ 
tako da je 
\begin{align}\label{t3.1}
    n>k_{0}(\varepsilon) \;\Rightarrow\; |a_{n}-a|<\varepsilon, 
    \;\;\;\;\; 
    n>l_{0}(\varepsilon) \;\Rightarrow\; |b_n-b|<\varepsilon. 
\end{align}

(a) Ovaj deo teorem je trivijalno ispunje ako je $\alpha=\beta=0$. 
U nastavku dokaza, pretpostavljavamo da je $\alpha\neq 0$ ili 
$\beta\neq 0$. 
Neka je $\gamma=|\alpha|+|\beta|$. 
Kako je $\alpha\neq 0$ ili $\beta\neq 0$, 
imamo $\gamma>0$. 
Pored toga, koristeћi (\ref{t3.1}), dobijamo 
\begin{align}\label{t3.3}
    |(\alpha a_n + \beta b_n) - (\alpha a + \beta b)|
    =&
    |\alpha(a_n-a)+\beta(b_n-b)|
    \nonumber\\
    \leq &
    |\alpha|\:|a_n-a| 
    + 
    |\beta|\:|b_n-b|
    \nonumber\\
    \leq &
    (|\alpha|+|\beta|)
    \max\{|a_n-a|,|b_n-b| \}
    \nonumber\\
    \leq &
    \gamma\varepsilon
\end{align}
za $n\geq\max\{k_{0}(\varepsilon),l_{0}(\varepsilon)\}$. 
Neka je 
$n_{0}(\varepsilon)=\max\{k_{0}(\varepsilon/\gamma),
l_{0}(\varepsilon/\gamma)\}$, gde je $\varepsilon$ 
pozitivan realan broj. 
Onda, zamenjujuћi $\varepsilon$ u (\ref{t3.3})  
vrednoshћu $\varepsilon/\gamma$, dobijamo 
\begin{align*}
    |(\alpha a_n + \beta b_n) - (\alpha a + \beta b)|<\varepsilon
\end{align*}
za $n>n_{0}(\varepsilon)$. 
Koristeћi definiciju granichne vrednosti i konvergentong niza, 
zakljuchujemo onda da je tvrdjenje dela (a) tachno. 

(b) Na osnovu dela (a) teoreme \ref{thm2}, 
postoji pozitivan realan broj $c$ tako da je 
$|a_n|\leq c$ za $n\in\mathbb{N}$. 
Neka je $d=|b|+c$. 
Kako je $c>0$, imamo $d>0$. 
Pored toga, koristeћi (\ref{t3.1}), zakljuchujemo  
\begin{align}\label{t3.5}
    |a_n b_n - ab|
    =&
    |a_n(b_n-b)+b(a_n-a)|
    \nonumber\\
    \leq &
    |a_n|\:|b_n-b| 
    + 
    |b|\:|a_n-a|
    \nonumber\\
    \leq &
    (|b|+d)
    \max\{|a_n-a|,|b_n-b| \}
    \nonumber\\
    \leq &
    d\varepsilon
\end{align}
za $n\geq\max\{k_{0}(\varepsilon),l_{0}(\varepsilon)\}$. 
Neka je 
$n_{0}(\varepsilon)=\max\{k_{0}(\varepsilon/d),
l_{0}(\varepsilon/d)\}$, gde je $\varepsilon$ 
pozitivan realan broj. 
Onda, zamenjujuћi $\varepsilon$ u (\ref{t3.5})  
vrednoshћu $\varepsilon/d$, dobijamo 
$|a_n b_n - ab|<\varepsilon$ 
za $n>n_{0}(\varepsilon)$. 
Koristeћi definiciju granichne vrednosti i konvergentong niza, 
zakljuchujemo onda da je tvrdjenje dela (b) tachno. 

(v) Kako je $b\neq 0$ i $b_n\neq 0$, imamo 
da je $|b|>0$ i $|b_n|>0$ za svako $n\in\mathbb{N}$.
Neka je $m=l_0(|b|/2)$ i neka je 
\begin{align*}
    \alpha=\min\{|b_1|,\dots,|b_m|\}, 
    \;\;\;\;\; 
    \beta=\min\left\{\alpha,|b|/2\right\}. 
\end{align*}
Onda, koristeћi (\ref{t3.1}), dobijamo 
\begin{align}\label{t3.7}
    |b_n|=|(b_n-b)+b|
    \geq 
    |b|-|b_n-b|
    \geq 
    \frac{|b|}{2}
    \geq 
    \beta
\end{align}
za $n>m=l_0(|b|/2)$. 
Pored toga, imamo  
\begin{align}\label{t3.9}
    |b_n|
    \geq 
    \min\{|b_1|,\dots,|b_m|\}
    =
    \alpha
    \geq \beta
\end{align}
za $1\leq n\leq m$. 
Iz (\ref{t3.7}) i (\ref{t3.9}) sledi 
$|b_n|\geq\beta$ za $n\in\mathbb{N}$. 

Neka je $c=(1+|a|/|b|)/\beta$. 
Kako je $|b|>0$ i $|b_n|>0$ za svako $n\in\mathbb{N}$, 
imamo da je $\beta>0$ i $c>0$.  
Pored toga, koristeћi (\ref{t3.1}), dobijamo  
\begin{align}\label{t3.21}
    \left|\frac{a_n}{b_n} - \frac{a}{b}\right|
    =&
    \left|\frac{b(a_n-a)-a(b_n-b)}{b\:b_n} \right|
    \nonumber\\
    \leq &
    \frac{|b|\: |a_n-a| + |a|\: |b_n-b|}{|b|\: |b_n|}
    \nonumber\\
    \leq &
   \frac{|a|+|b|}{\beta |b|}
    \max\{|a_n-a|,|b_n-b| \}
    \nonumber\\
    \leq &
    c\varepsilon
\end{align}
za $n\geq\max\{k_{0}(\varepsilon),l_{0}(\varepsilon)\}$. 
Neka je 
$n_{0}(\varepsilon)=\max\{k_{0}(\varepsilon/c),
l_{0}(\varepsilon/c)\}$, gde je $\varepsilon$ 
pozitivan realan broj. 
Onda, zamenjujuћi $\varepsilon$ u (\ref{t3.21})  
vrednoshћu $\varepsilon/c$, dobijamo 
\begin{align*}
    \left|\frac{a_n}{b_n} - \frac{a}{b}\right|<\varepsilon
\end{align*}
za $n>n_{0}(\varepsilon)$. 
Koristeћi definiciju granichne vrednosti i konvergentong niza, 
zakljuchujemo onda da je tvrdjenje dela (c) tachno. 
\end{dokaz}

Teorema \ref{thm3} ilustrovana je sledeћim primerima. 

\begin{pri}
(a) Neka je $(a_n)_{n\in\mathbb{N}}$ harmonijski niz 
(tj. $a_n=1/n$ za svako $n\in\mathbb{N}$). 
Pokazhimo da je $(a_n)_{n\in\mathbb{N}}$ nula-niz. 
Neka je $(b_n)_{n\in\mathbb{N}}$ niz chiji je $n$-ti
chlan $b_n=(n-1)/n$. 
Onda imamo $a_n=1-b_n$, 
dok deo (a) primera \ref{pri3.1} daje 
$\displaystyle\lim_{n\rightarrow\infty}b_n=1$. 
Koristeћi deo (a) teoreme \ref{thm3} dobijamo 
\begin{align*}
    \lim_{n\rightarrow\infty}
    a_n=
    \lim_{n\rightarrow\infty}(1-b_n)
    =
    1-\lim_{n\rightarrow\infty}b_n
    =0, 
\end{align*}
tj. $(a_n)_{n\in\mathbb{N}}$ je nula-niz. 

(b) Neka je $(a_n)_{n\in\mathbb{N}}$ niz chiji je $n$-ti 
chlan 
\begin{align*}
    a_n=\frac{3-2n+5}{4n^2+n-6}. 
\end{align*}
Pokazhimo da je $\displaystyle\lim_{n\rightarrow\infty}a_n=3/4$. 
Neka su $(b_n)_{n\in\mathbb{N}}$ i $(c_n)_{n\in\mathbb{N}}$ 
nizovi ciji su 
$n$-ti chlaovi 
\begin{align*}
    b_n=3-\frac{2}{n}+\frac{5}{n^2}, 
    \;\;\;\;\; 
     c_n=4+\frac{1}{n}-\frac{6}{n^2}. 
\end{align*}
Onda imamo $a_n=b_n/c_n$, 
dok delovi (a), (b) Teoreme \ref{thm3} i deo (a) ovog primera daju 
\begin{align}
    &\label{p4.1.1}
    \lim_{n\rightarrow\infty} b_n
    =
    \lim_{n\rightarrow\infty}
    \left(3-\frac{2}{n}+\frac{5}{n^2} \right)
    =
    3 
    - 
    2\lim_{n\rightarrow\infty}\frac{1}{n}
    +
    5\left(\lim_{n\rightarrow\infty}\frac{1}{n} \right)^2
    =
    3, 
    \\
    &\label{p4.1.3}
    \lim_{n\rightarrow\infty} c_n
    =
    \lim_{n\rightarrow\infty}
    \left(4+\frac{1}{n}-\frac{6}{n^2} \right)
    =
    4 
    + 
    \lim_{n\rightarrow\infty}\frac{1}{n}
    -
    6\left(\lim_{n\rightarrow\infty}\frac{1}{n} \right)^2
    =
    4. 
\end{align}
Koristeћi deo (b) teoreme \ref{thm3} i (\ref{p4.1.1}), (\ref{p4.1.3}), 
dobijamo 
\begin{align*}
    \lim_{n\rightarrow\infty}
    a_n
    =
    \lim_{n\rightarrow\infty}\frac{b_n}{c_n}
    =
    \frac{\displaystyle\lim_{n\rightarrow\infty}b_n}
    {\displaystyle\lim_{n\rightarrow\infty}c_n}
    =\frac{3}{4}.  
\end{align*}
\end{pri}

\begin{thm}\label{thm4}
Neka su $(a_n)_{n\in\mathbb{N}}$, 
$(b_n)_{n\in\mathbb{N}}$ i 
$(c_n)_{n\in\mathbb{N}}$ nizovi realnih brojeva 
sa sledeћim osobinama: 

(a) $a_n\leq c_n \leq b_n$ za svako $n\in\mathbb{N}$. 

(b) $(a_n)_{n\in\mathbb{N}}$ i $(b_n)_{n\in\mathbb{N}}$
su konvergentni i ispunjavaju 
$\displaystyle\lim_{n\rightarrow\infty} a_n = 
\displaystyle\lim_{n\rightarrow\infty} b_n$. 
\\
Onda je $(c_n)_{n\in\mathbb{N}}$ konvergentan niz za koji vazhi $\displaystyle\lim_{n\rightarrow\infty} c_n =
\displaystyle\lim_{n\rightarrow\infty} a_n$. 
\end{thm}

\begin{nap}
U anglosaksonskoj literaturi, 
Teorema \ref{thm4} naziva se sendvich teorema. 
Ova teorema pokazuje da je niz konvergentan ako je ,,ukljeshten''
izmedju dva konvergentna niza sa istim granichnim vrednostima. 
\end{nap}

\begin{dokaz}
Videti referencu \cite{kadelburg}. 
\end{dokaz}

\begin{thm}\label{thm5}
(a) Neka je $(a_n)_{n\in\mathbb{N}}$ monotono rastuћi i 
odozgo ogranichen niz realnih brojeva. 
Onda je $(a_n)_{n\in\mathbb{N}}$ konvergentan. 

(b) Neka je $(b_n)_{n\in\mathbb{N}}$ monotono opadajuћi i 
odozdo ogranichen niz realnih brojeva. 
Onda je $(b_n)_{n\in\mathbb{N}}$ konvergentan. 
\end{thm}

\begin{nap}
Teorema \ref{thm5} uspostavlja vezu izmedju 
konvergenicije, ogranichenosti i monotonosti nizova. 
Na osnovu ove teoreme, monoton i ogranichen niz je uvek konvergentan. 
\end{nap}

\begin{dokaz}
Videti referencu \cite{kadelburg}. 
\end{dokaz}

Teoreme \ref{thm4} i \ref{thm5} ilustrovane su sledeћim primerima. 

\begin{pri}
(a) Neka je $q$ realan broj takav da je $0\leq q<1$
i neka je $(a_n)_{n\in\mathbb{N}}$ niz chiji je $n$-ti chan 
$a_n=q^n$. 
Pokazhimo da je $\displaystyle\lim_{n\rightarrow\infty}a_n=0$. 
Kako je $a_{n+1}=a_n q$ i $0\leq q<1$, 
$(a_n)_{n\in\mathbb{N}}$ je monotono opadajuћi niz. 
S obzirom da je $a_n\geq 0$, iz dela (b) teoreme \ref{thm5} 
sledi da je $(a_n)_{n\in\mathbb{N}}$ konvergentan, 
tj. postoji realan broj $a$ tako da je 
$\displaystyle\lim_{n\rightarrow\infty}a_n=a$. 
Kako je $a_{n+1}=a_n q$, deo (a) teorem \ref{thm3} daje 
\begin{align*}
    a
    =
    \lim_{n\rightarrow\infty}a_{n+1}
    =
    \lim_{n\rightarrow\infty}(a_{n}q)
    =
    q\lim_{n\rightarrow\infty}a_{n}
    =
    aq, 
\end{align*}
tj. $a(q-1)=0$. 
S obzirom da je $q\neq 1$, 
zakljuchujemo da je $a=0$, 
tj. $\displaystyle\lim_{n\rightarrow\infty}a_=0$. 

(b) Neka su $a$ i $q$ realni brojevi takvi da je $-1<q<1$ 
i neka je neka je $(a_n)_{n\in\mathbb{N}}$ harmonijski 
niz chiji je $n$-ti chan 
$a_n=aq^n$. 
Pokazhimo da je $\displaystyle\lim_{n\rightarrow\infty}a_n=0$. 
Neka su $(b_n)_{n\in\mathbb{N}}$ i $(c_n)_{n\in\mathbb{N}}$
nizovi chiji su $n$-ti chlanovi 
$b_n=-|a|\:|q|^n$ i $c_n=|a|\:|q|^n$. 
Kako je $0\leq |q|<1$, deo (a) teorem \ref{thm3} 
i deo (a) ovog primera daju 
\begin{align}
    &\label{p4.2.1}
    \lim_{n\rightarrow\infty}b_n
    =
    \lim_{n\rightarrow\infty}(-|a|\:|q|^n)
    =
    -|a|\lim_{n\rightarrow\infty}|q|^n
    =
    0, 
    \\
    &\label{p4.2.3}
    \lim_{n\rightarrow\infty}c_n
    =
    \lim_{n\rightarrow\infty}(|a|\:|q|^n)
    =
    |a|\lim_{n\rightarrow\infty}|q|^n
    =
    0. 
\end{align}
S obzirom da je $b_n\leq a_n\leq c_n$, 
iz teoreme \ref{thm4} i (\ref{p4.2.1}), (\ref{p4.2.3})
sledi da je 
$\displaystyle\lim_{n\rightarrow\infty}a_n=
\displaystyle\lim_{n\rightarrow\infty}b_n=0$. 
\end{pri}

\begin{thm}\label{thm6}
(a) Neka je $(a_n)_{n\in\mathbb{N}}$ odredjeno divergentan niz 
realnih brojeva takav da je $a_n\neq 0$ za svako $n\in\mathbb{N}$.
Onda je $\displaystyle\lim_{n\rightarrow\infty}1/a_n=0$
(tj. $1/(a_n)_{n\in\mathbb{N}}$ je nula niz). 

(b) Neka je $(b_n)_{n\in\mathbb{N}}$ nula niz 
realnih brojeva takav da je $b_n\neq 0$ za svako $n\in\mathbb{N}$. 
Ako postoji prirodan broj $k$ tako da je $b_n>0$ za $n>k$, 
onda je $\displaystyle\lim_{n\rightarrow\infty}1/b_n=\infty$
(tj. $1/(b_n)_{n\in\mathbb{N}}$ je odredjeno divergentan niz).

(v) Neka je $(c_n)_{n\in\mathbb{N}}$ nula niz 
realnih brojeva takav da je $c_n\neq 0$ za svako $n\in\mathbb{N}$. 
Ako postoji prirodan broj $l$ tako da je $c_n<0$ za $n>l$, 
onda je $\displaystyle\lim_{n\rightarrow\infty}1/c_n=-\infty$
(tj. $1/(c_n)_{n\in\mathbb{N}}$ je odredjeno divergentan niz).
\end{thm}

\begin{nap}
Teorema \ref{thm6} uspostavlja vezu izmedju odredjeno divergentnih 
i nula nizova. 
Na osnovu dela (a), 
niz reciprochan odredjeno divergentnom nizu je nula-niz. 
Deo (b) pokazuje da je niz 
reciprochan nula-nizu pozitivnih realnih brojeva odredjeno divergentan. 
Na osnovu dela (v), isto vazhi za niz reciprochan 
nula-nizu negativnih realnih brojeva. 
\end{nap}

\begin{dokaz}
Videti referencu \cite{kadelburg}. 
\end{dokaz}

\begin{thebibliography}{10}

\bibitem{green}
\lat
J. A. Green, {\em Sequences and Series,}
{Rountledge} \& Kegan Paul Ltd 
and Dover Publications, New York, 1958.  
\cyr
\bibitem{kadelburg}
Zoran Kadelburg, Vladimir Miћић, Srdjan Ognjanoviћ 
i Sonja Чукић,  
{\em Analiza sa Algebrom, Udџbenik sa zbirkom zadataka za 3. razred Matematiчke gimnazije,}
Krug, Beograd, 2018. 

\end{thebibliography}

\end{document}
